\subsection{Observed tissue composition provides the context for a fibroblast readout}

The study material comprises dissociated foot-skin specimens from diabetic-ulcer and non-diabetic groups, rather than longitudinally tracked cells (Figure~\ref{fig:p1f1}A\label{call:p1f1}). The 25 specimens represent 7 healed, 4 non-healed and 9 non-diabetic patients; cohort roles and analysis units are specified in Table~\ref{tab:p1cohorts}. The schematic in A separates model training from frozen inference, then branches into the fibroblast-restricted mean/fraction and the primary all-cell patient contrast. These are different estimands, not successive processing stages. The encoder and decoder in Equations~\eqref{eq:p1encoder} and~\eqref{eq:p1decoder} define the measured coordinate; the centered-log-ratio display in Equation~\eqref{eq:p1display} changes only its visualization. Mapping the saved expression coordinates makes the recovered cellular populations visible before comparing their statistical summaries. The display includes stromal, epithelial, vascular and immune compartments among 76,461 cells; fibroblasts account for 19,410 cells, or 25.4\% of the retained population (Figure~\ref{fig:p1f1}B). Coloring the same coordinates by the frozen fibroblast-associated loading shows how that continuous measurement relates to the operational lineage assignments (Figure~\ref{fig:p1f1}C). No new cell type or trajectory is defined by separation in this two-dimensional display.
\paperfigure{figure1_workflow}{fig:p1f1}{Observed study material and expression context. (A) Conceptual skin anatomy and sampling, not a tissue image or inferred spatial map. (B) A uniform 2,000-cell subset of the 76,461-cell UMAP, colored by marker-derived lineage. Fibro., fibroblast; EC, endothelial cell; mural, pericyte/smooth muscle. (C) The identical cells colored by frozen programme loading; the map and colors do not define new cell types. (D) Raw expression across all retained cells: dot size is the detection fraction and color the gene-wise scaled mean log-normalized expression. Annotation markers provide descriptive, not independent, support for labels. (E) All 19,410 fibroblast loadings in fixed-width histograms, with the existing threshold marked; healthy participants have no healing outcome. (F) Recovered-cell composition in each of 25 specimens, grouped by source outcome or non-diabetic status. Colors match B; columns are specimens, including repeated patients. Inference uses the relevant full data and patient units, not the display subsample.}


Raw-count reconstruction reproduced all saved lineage assignments and their sample order. The measured marker profiles provide an expression-level view of those assignments and of the programme-associated genes (Figure~\ref{fig:p1f1}D). Because some displayed markers also contributed to annotation, this is descriptive support rather than an independent validation of cell identity. Loading distributions within fibroblasts retain both low and high values across the clinical groups (Figure~\ref{fig:p1f1}E), while individual specimens differ in their recovered cell composition (Figure~\ref{fig:p1f1}F). Consequently, a pooled expression map, a cell fraction and a patient outcome contrast describe different aspects of the material. The mapped patient readouts are retained in Table~\ref{tab:p1biologicalpatients}.
\FloatBarrier

\subsection{The frozen representation is structured, but its principal loading is a mixture coordinate}
The discovery representation retained 76,461 quality-controlled cells from 25 selected foot-skin samples in the original 54-sample series. Beta perplexity was 243, and the manifold audit passed its recorded integrity checks. The connectivity statistic in Equation~\eqref{eq:p1reach} does not establish local geometric fidelity. The fitted loss in Equation~\eqref{eq:p1loss} uses Gaussian sampling and regularization as specified in Equation~\eqref{eq:p1gaussian}; deterministic readouts use the encoder mean, not a newly sampled latent state. Fibroblast topic was enriched for TIMP1, COL3A1, DCN, FN1, LUM and CHI3L1, with immediate-early and oxidative-stress genes also present among the top-ranked terms. The gene list therefore supports a fibroblast-associated inflammatory/stress description; it does not by itself establish a unique cell state or a healing mechanism.

Applying the mean--fraction identity in Equation~\eqref{eq:p1identity} to all 19,410 discovery fibroblasts clarified the remaining within-bin contribution. Empirical low-bin means ranged from 0.0120 to 0.1021, and nonempty high-bin means from 0.2264 to 0.5778. Nine of 25 specimens had an empty high bin. The approximation residual defined in Equation~\eqref{eq:p1residual} had RMS of 0.01575 and a residual range of $-0.02434$ to $0.04377$. The exact mean identity held to below $10^{-15}$ in every specimen (Table~\ref{tab:p1math}). The constant-state approximation is therefore explicit and imperfect; the correlation cannot establish invariant programme intensity.

A historical exploratory scan of 31 covariates across 25 specimens yielded 14 numerical BH q values below 0.05 under a naive independent-row analysis. Repeated patients invalidate that row-independence premise, so these are not 14 false-discovery-rate-controlled findings. The largest positive association was with the assigned B/plasma-cell fraction ($\rho=0.828$); CFD/APOD-rich topic mean loading was negatively associated ($\rho=-0.766$). Fibroblast PTPRC positivity and median UMI count were also associated ($\rho=0.727$ and 0.663). The complete scan, including non-significant covariates, is in Table~\ref{tab:p1covariates}. These marginal, sample-unit correlations do not establish causal drivers or independence from tissue composition; the positive depth association remains relevant despite the successful count-matching control.

\Needspace{65pt}
The hard-threshold fraction in Equation~\eqref{eq:p1fraction} uses a pooled two-Gaussian fit for a descriptive threshold; this fit did not test a continuous-shift model against a biological mixture model. High-loading samples retained a low-state tail overlapping the low-loading samples. Across the 25 analyzable samples, the sample mean and hard-threshold high-state fraction had Spearman $\rho=0.933$ (Figure~\ref{fig:p1f2}A\label{call:p1f2}). Quantile-based sample categories summarize the coexistence of low and high loadings without assigning a biological mechanism to those categories (Figure~\ref{fig:p1f2}B). The historical specimen-bootstrap interval of 0.789--0.983 and specimen-omission range of 0.924--0.957 are retained only as row-level diagnostics in Table~\ref{tab:p1semantics}, not patient-level uncertainty. The permutation null was centred near zero (Table~\ref{tab:p1semantics}). This is the central measurement result: in this dataset, the sample average is strongly associated with the high-state fraction. The two summaries are calculated from the same cells, so their correlation alone does not prove that component intensities are constant. The result does not say that high-state cells have no intrinsic biology. Together with the low-state tail, it shows that the mean contains information about high-score prevalence; it does not identify the relative biological contributions of prevalence and within-state expression. Patient aggregation gave a correlation of 0.937 across 20 patients, with a 10,000-draw patient-bootstrap 95\% interval of 0.776--0.990. Omitting one patient at a time gave correlations of 0.927--0.964 (Figure~\ref{fig:p1f2}C; Table~\ref{tab:p1robustness}). These descriptive patient-resampling summaries replace specimen resampling as the primary uncertainty display, without adding independent observations.

\paperfigure{figure2_mixture_semantics}{fig:p1f2}{Association between mean topic loading within fibroblasts and the high-state fraction. (A) Each point is one of 25 GSE165816 specimens: 9 healing (green circles), 5 non-healing (orange triangles) and 11 healthy non-diabetic (purple squares). The high-state fraction uses the pooled cut 0.1842958; Spearman ρ=0.933. (B) Sample-shape counts use within-sample loading quantiles: pure low requires the 90th percentile at most 0.15, pure high the 10th percentile at least 0.10, and mixed the 10th percentile at most 0.05 and 90th percentile at least 0.40; other samples form the remaining category. (C) Unlike the specimen-level panels A and B, uncertainty uses 20 author-mapped patients after pooling repeated specimens by fibroblast count. The dot is the patient correlation 0.937, the blue segment is the descriptive 10,000-draw patient-bootstrap 95\% interval (0.776--0.990), and the orange segment is the leave-one-patient-out range (0.927--0.964). Both correlated summaries use the same cells; the association does not by itself prove stable within-state intensity or a healing effect.}


\subsection{The operational high-state pattern survives selected technical controls}

The operational high-state pattern survived the originally specified artifact challenges, while the expanded controls identify measurement dependence. The dissociation stress standardized mean difference was +0.246, below the rejection threshold. PTPRC positivity was 0.98\% in high-state cells and 0.62\% in low-state cells, a rate ratio of 1.60 (Figure~\ref{fig:p1f3}A\label{call:p1f3}). Exact panel-count matching produced 99.5\% cell-label agreement and a specimen high-state-fraction correlation of 0.984 among 14,561 higher-depth retained fibroblasts under fixed lineage labels (Figure~\ref{fig:p1f3}B). This does not test the 4,849 excluded fibroblasts or rerun the quality-control and lineage-assignment pipeline after count reduction. The result therefore supports conditional score stability, not depth invariance of the entire tissue measurement. PTPRC positivity is only a proxy and does not exclude all doublets. A raw-count structureless control preserved every gene\textquotesingle{}s value multiset but changed the topic concentration substantially (perplexity 3.44 to 11.23 of 15), showing sensitivity to the joint perturbation of co-expression and row-wise library-size and detection-rate variation. The adaptive marker reference changed 9.21\% of cell labels and shifted the maximum specimen mean by 0.0371 and high-state fraction by 0.0695; there is no independent identity truth set to select between gates (Table~\ref{tab:p1gatecontrol}).

\paperfigure{figure3_artifact_controls}{fig:p1f3}{Technical and immune-transcript exclusion controls in the discovery cohort. (A) Dissociation standardized mean difference (SMD, 0.246) and the high-state/low-state PTPRC-positive rate ratio (1.60) are divided by rejection bounds 0.8 and 2.0; the dashed line is one. The rounded state cut is 0.184. (B) Each point is one of 25 samples after restricting to 14,561 fibroblasts eligible for panel-count matching to 986 counts per cell. Reprojection preserves sample fractions (ρ=0.984) and 99.5\% of cell labels; the dashed line is identity. (C) Bars are BIC for one Gaussian minus BIC for two Gaussians fitted to sample high-state fractions, using 24 matched samples before and after exclusion. Positive values favour two components. Exclusion retains 11,262 of 19,410 fibroblasts overall, and matched sample fractions correlate at ρ=0.989. This control measures invariance to the exclusion rule; it does not identify every source of ambient RNA or tissue-composition dependence.}


The full released-count audit reproduced the 81,602 pre-QC cells, 76,461 baseline survivors and all 19,410 fibroblast assignments, with a maximum deterministic-coordinate discrepancy below $4.77\times10^{-7}$. At 75\%, 50\%, 25\% and 10\% count retention, the three seeded runs retained 76,375--76,393, 75,724--75,738, 72,145--72,197 and 61,366--61,397 cells, respectively. Among cells surviving both QC passes, 3.18--3.30\%, 6.23--6.34\%, 10.29--10.48\% and 15.31--15.40\% changed lineage labels. At 10\%, only 80.11--80.15\% of original survivors remained, while 111--115 original failures newly passed QC. The gate assigned a winning label to 92 baseline QC survivors with no detected annotation marker; this increased to 1,410--1,462 at 10\% depth (Table~\ref{tab:p1fulldepth}). These labels are not identity truth. At 10\%, score rank correlation was about 0.910 on original retained cells with estimable scores, versus about 0.947 among QC survivors, exposing selection in a survivor-only stability claim. The full-pipeline all-cell patient contrast ranged from 0.0288 to 0.0296, with exact probabilities 109/330--114/330, versus baseline 90/330. All four readouts and rates are retained in Table~\ref{tab:p1depthoutcomes}. Persistence of a group contrast therefore coexists with substantial changes to the cells and labels being summarized.

Removing cells carrying any of the prespecified immune transcripts left the specimen high-state fractions nearly unchanged ($\rho=0.989$) and increased the BIC support for two components fitted to the sample fractions (Figure~\ref{fig:p1f3}C). This supports stability under the specified immune-transcript exclusion rule; it does not identify or exclude ambient RNA contamination. The placebo analysis was deliberately retained: unrelated tissue covariates flattened residual bimodality as well as immune covariates. At 24 matched samples, the data cannot determine whether the state is statistically separable from every tissue-composition axis. The matched exclusion comparison contains 24 samples, while each residual regression uses 25 samples. The immune and placebo regressions explained 49.5\% and 35.5\% of the variance, respectively, and both residual BIC differences favoured one component (Table~\ref{tab:p1ambient}). A clean technical control is not evidence that the state is a healing marker.

\subsection{High-state abundance does not establish healing specificity}

The lineage marker list shared no genes with the specified top-gene list used to describe the tested topic. This limited non-overlap does not make the gate independent of the full-panel learned representation. Fibroblast composition itself differed between the healing and non-healing sample groups only descriptively (ratio 1.42, $p=0.438$). Within the fibroblast gate, the discovery contrast was 2.89-fold, with a bootstrap 95\% interval of 1.06--13.97 and sample-level $p=0.04196$. This value is retained as a discovery observation, not as an independent patient validation, because the 14 samples represent 11 patients and include repeated same-wound, same-day specimens. The complete lineage and composition contrasts are in Table~\ref{tab:p1lineage}.

The high-state pattern was not restricted to samples labelled as healing. Four of nine healing samples and none of five non-healing samples were classified as high-mode samples, Fisher exact $p=0.2208$. The most extreme high-state fraction was a healthy non-diabetic foot-skin sample, whose high-state fraction was 0.841 (Figure~\ref{fig:p1f4}A\label{call:p1f4}). That observation is an existence counterexample to the statement that the state is uniquely a marker of a healing DFU. It neither identifies the tissue process responsible nor refutes a possible probabilistic association with healing.

\paperfigure{figure4_outcome_counterexample}{fig:p1f4}{State abundance, analysis unit and design power. (A) Symbols show high-state fractions for 9 healing, 5 non-healing and 11 healthy specimens; horizontal segments mark arm means. Green circles, orange triangles and purple squares identify these groups. The largest fraction is 0.841 in healthy non-diabetic foot skin, showing that the state is not exclusive to healing ulcers. (B) Dots and within-arm bootstrap 95\% intervals show the healing-minus-non-healing difference in all-cell mean fibroblast-associated topic loading for 14 specimens and 11 mapped patients. The patient analysis resamples 7 healing and 4 non-healing patients in 30,000 draws. Its readout differs from the fibroblast high-state fraction in A; the vertical dashed line is zero. (C) Gaussian location-shift simulation with a two-sided Mann--Whitney test estimates 24.9\% power for the separate 9-versus-5 high-state-fraction design, against an 80\% target. Its minimum detectable difference is 0.519. These are design sensitivities, not effect-exclusion bounds or clinical minimum important differences.}


For the separate all-cell mean-loading contrast in Equation~\eqref{eq:p1delta}, collapsing repeats to 11 patients gave a healing-minus-non-healing difference of 0.0293 (Figure~\ref{fig:p1f4}B). This is not a patient-level recalculation of the 2.89-fold fibroblast-restricted contrast. Exhaustive enumeration in Equation~\eqref{eq:p1exact} gave a two-sided probability of 0.272727 (90/330 allocations at least as extreme), versus 0.274924 under the retained add-one rule in Equation~\eqref{eq:p1permutation}; the bootstrap 95\% interval was {[}−0.0024, 0.0667{]}. The Welch standard error was 0.019446, and the data-derived equivalence infimum was 0.065375. The earlier 14-sample diagnostic contrast was 0.0457 with exhaustive $p=0.078921$ and an equivalence infimum of 0.079644; it is retained as a specimen-unit sensitivity analysis. The two-one-sided-test rule in Equation~\eqref{eq:p1tost} and boundary in Equation~\eqref{eq:p1margin} explain these data-derived sensitivity limits; neither margin is a clinical MCID. The corrected result is therefore a unit-aware discovery result with substantial uncertainty, not a definitive null test (Table~\ref{tab:p1exact}). Mapped scores are given in Table~\ref{tab:p1patients}; equal weighting of repeated specimen means gave 0.0297. Restricting the secondary specimen-label null to the 14 DFU specimens, rather than allowing healthy controls into pseudo-outcome groups, gave exact probabilities 0.0894, 0.1199 and 0.2697 for the fibroblast mean, high-state fraction and fibroblast composition, respectively; repeated specimens make these secondary quantities (Table~\ref{tab:p1dfunull}). A constructed continuous-score example produced a median mean--fraction Spearman correlation of 0.951, at least as large as the observed 0.933 in 76\% of draws, so the correlation cannot identify a discrete mixture without further assumptions. This construction is illustrative rather than a fitted null or calibrated probability. The full equivalence sensitivity grid is in Table~\ref{tab:p1equivalence}. Patient omission left the all-cell contrast positive, from 0.0143 to 0.0365 across 11 deletions (Table~\ref{tab:p1robustness}); this is an influence range, not a confidence interval or evidence of clinical discrimination.


\subsection{Retrospective design simulations do not resolve outcome uncertainty}

A retrospective simulation substituted the observed specimen-level high-state-fraction effect and pooled spread into a hypothetical 9-versus-5 independent-specimen design, giving 24.9\% estimated power. Under that simulation, the difference associated with 80\% power was 0.519 (Figure~\ref{fig:p1f4}C). The observed specimens include repeated patients, so this is neither the power of the corrected 7-versus-4 patient comparison nor independent evidence explaining its nonsignificance. These conditional design calculations cannot exclude a moderate or large healing association.

We also ran a design calculation for the next cohort, keeping it separate from the biological results. Under the deterministic specimen-derived arm spreads, hypothetical independent observations, 80\% target power and no missing outcomes, a candidate difference of 0.03 gave a total of 47, while a candidate equivalence margin of 0.05 gave 20. These spreads include repeated specimens and are not patient-population variance estimates. These values are design inputs rather than findings: the clinical MCID is still open, repeated samples increase the design effect, and the calculation cannot turn the discovery cohort into an adequately powered study. Tables~\ref{tab:p1design} and~\ref{tab:p1aucdesign} give the full difference, equivalence and exploratory AUC-gap design calculations.

\subsection{Fibroblast-restricted patient comparisons retain substantial uncertainty}

We next evaluated the two fibroblast-restricted readouts at the same independent patient unit used for the all-cell analysis. This closes the gap between the original sample-level lineage comparison and the mapped clinical outcomes. The denominator-specific pooling rules in Equation~\eqref{eq:p1lineagepatients} gave mean within-fibroblast loadings of 0.1729 in seven healed patients and 0.0837 in four non-healed patients. Their difference was 0.0892, with a 10,000-draw patient-bootstrap interval of $[-0.0486,0.2177]$ and an exhaustive two-sided permutation probability of 0.3121 (Figure~\ref{fig:p1f5}A\label{call:p1f5}). The corresponding high-state fractions were 0.3010 and 0.1496; the difference was 0.1514, with interval $[-0.1331,0.4153]$ and $p=0.3697$ (Figure~\ref{fig:p1f5}B). These comparisons use the existing gate and threshold, without searching for a cut that improves outcome separation.

\paperfigure{figure5_sensitivity_negative_controls}{fig:p1f5}{Threshold, capacity and negative-control analyses. (A,B) Spearman correlation of fibroblast mean loading with the high-state fraction across 25 samples and the two-sided Mann--Whitney healing-arm p value across 9 versus 5 samples, for all 18 cuts from 0.08 to 0.40. Vertical dashed lines mark the fitted cut; the horizontal line in B marks p=0.05. (C) Correlations for K=10, 15 and 20 after tracking the TIMP1-carrying topic; orange marks K=10, where the highest fraction is no longer in healthy skin. (D) The observed high-state-fraction arm difference is compared with the central 95\% of 20,000 label permutations (p=0.1358). (E) Gene-column shuffling of 6,000 cells raises mean cell perplexity from 3.68 to 14.43 out of 15, violating the concentration criterion. (F) The anatomical sample-level control compares 33 foot and 11 forearm specimens: observed p=0.0015, with 5.0\% of permuted comparisons below 0.05. This calibration is distinct from the paired 10-patient anatomy analysis and does not establish pipeline-wide error control.}


Fibroblast composition was evaluated separately, retaining all recovered cells as its denominator. The patient-level fractions were 0.2825 and 0.2118, a difference of 0.0708 with interval $[-0.0708,0.2146]$ and $p=0.4697$ (Figure~\ref{fig:p1f5}C). Benjamini--Hochberg adjustment across these three new readouts gave $q=0.4697$ for each. These are retrospective within-cohort comparisons; neither the positive point estimates nor their lack of significance establishes prognostic validity or equivalence.

Equal weighting of repeated specimen summaries within each patient gave differences of 0.0974, 0.1683 and 0.0679 for the programme mean, high-state fraction and fibroblast composition, respectively. All DFU specimens already contained at least 100 assigned fibroblasts, so imposing that requirement changed no outcome observation and supplies no independent robustness evidence. Omitting one patient at a time left the three differences positive, with ranges 0.0593--0.1445, 0.0850--0.2720 and 0.0279--0.1173, respectively. These influence ranges coexist with bootstrap intervals spanning zero; they are not confidence bounds. Table~\ref{tab:p1biologicalcontrasts} reports the complete weighting comparisons and exact probabilities.

The threshold-free patient distribution statistic in Equation~\eqref{eq:p1ecdf} was 0.3045, with exhaustive $p=0.1848$ over 330 patient-label allocations. Across the eight fixed cuts and the original midpoint, high-fraction differences remained positive (0.0570--0.2485); all unadjusted probabilities exceeded 0.22 and all support-wide max-statistic probabilities exceeded 0.31 (Figure~\ref{fig:p1f5}D; Table~\ref{tab:p1thresholdpatient}). Thus, the uncertain outcome association persists when the evaluation considers the whole loading distribution or multiple fixed thresholds, instead of relying only on the fitted midpoint.

Equalizing each specimen to 100 sampled fibroblasts left the programme-mean contrast positive in all 250 draws (0.0749--0.1233) and the high-state contrast positive in all draws (0.1098--0.2229); their smallest exact probabilities were 0.1273 and 0.1515 (Figure~\ref{fig:p1f5}E). Selecting one specimen per patient in all eight possible combinations gave programme-mean, high-state and composition ranges of 0.0777--0.1170, 0.1284--0.2081 and 0.0440--0.0918, respectively (Figure~\ref{fig:p1f5}F; Table~\ref{tab:p1samplingsensitivity}). These computations narrow the explanation that the observed direction is solely caused by unequal recovery or a particular repeated specimen, while retaining uncertainty about the healing association.


Direct comparison of the five repeated-specimen pairs, including three DFU patients, showed that within-specimen cell partitioning does not describe all between-specimen variation. In G7 and G8 from the same wound and day, fibroblast mean loading was 0.3504 versus 0.0921, an absolute difference of 0.2583. The separate within-specimen split-half 97.5th-percentile mean differences were 0.0553 and 0.0185; the full-distribution CDF supremum between specimens was 0.5789 (Table~\ref{tab:p1repeats}). This diagnostic leaves tissue microheterogeneity, preparation and annotation differences confounded. It does not estimate clinical repeatability or a patient-level significance probability.

\subsection{Method constants expose both stable and fragile claims}

The mean-versus-weight association remained positive across the cut grid, from $\rho=0.909$ to 0.948 (Figure~\ref{fig:p1f6}A\label{call:p1f6}), and across K=10, 15 and 20, from $\rho=0.924$ to 0.947 (Figure~\ref{fig:p1f6}C). The association therefore persisted over the tested cut and capacity choices.

\paperfigure{figure6_sensitivity_negative_controls}{fig:p1f6}{Threshold, capacity and negative-control analyses. (A) Each point gives the Spearman correlation between mean fibroblast loading and high-state fraction across 25 discovery specimens at one of 18 cuts from 0.08 to 0.40. The vertical dashed line marks the pooled fitted cut, 0.1842958. (B) Each point gives the two-sided Mann--Whitney healing-arm p value for 9 healing versus 5 non-healing specimens at the same cuts. The vertical line marks the fitted cut and the horizontal line marks p=0.05; 4 of 18 values fall below that level. (C) Correlations for K=10, 15 and 20 after tracking the TIMP1-carrying topic; orange marks K=10, where the highest fraction is no longer in healthy skin. (D) The observed high-state-fraction arm difference is compared with the central 95\% of 20,000 label permutations (p=0.1358). (E) Gene-column shuffling of 6,000 cells raises mean cell perplexity from 3.68 to 14.43 out of 15, violating the concentration criterion. (F) The anatomical sample-level control compares 33 foot and 11 forearm specimens: observed p=0.0015, with 5.0\% of permuted comparisons below 0.05. This calibration is distinct from the paired 10-patient anatomy analysis and does not establish pipeline-wide error control.}


The healthy-skin counterexample was more conditional. It held at all 18 tested cuts, but it failed at K=10 because the topic carrying TIMP1 belonged to a different IGFBP-rich programme. It held at K=15 and K=20 when the CHI3L1/COL3A1/DCN/FN1 state resolved as its own topic. The defensible wording is consequently ``the counterexample holds for the tested K=15 and K=20 fits but not the K=10 fit,'' not ``the counterexample is invariant to model capacity.''

The healing-arm p value moved between 0.042 and 0.062 over the cut grid, with only 4 of 18 cuts below 0.05 (Figure~\ref{fig:p1f6}B; Tables~\ref{tab:p1cuts} and~\ref{tab:p1capacity}). That sensitivity is a reason to avoid presenting the discovery p value as a stable healing signal. It also independently supports the decision to keep the healing claim unvalidated.

\subsection{Negative controls and representation comparisons constrain interpretation}

The DFU-only label-permutation calibration used the 9 healing and 5 non-healing specimens, excluding all 11 healthy controls before allocation. Its observed high-state-fraction difference was +0.2380, with a central 95\% null interval {[}−0.3167, +0.3222{]} and exact $p=0.1199$ over 2,002 allocations (Figure~\ref{fig:p1f6}D). Because this null relabels specimen outcomes and includes repeated patients, it is a secondary specimen-unit calibration rather than a patient-level clinical test. The observed contrast was not exceptional under this label-exchangeability null, and this single calibration does not prove pipeline-wide false-positive control. The raw-count structure control used a fixed 6,000-cell source-order subset and shuffled gene values across cells while preserving each gene\textquotesingle{}s observed value multiset; real data had perplexity 3.44/15, whereas the shuffled input had 11.23/15 and lower HHI (0.574 to 0.147). The shuffle therefore reduced coordinate concentration while also narrowing row-wise library-size and detection-rate variation; because it does not preserve row-wise library sizes, it is a co-expression-structure control rather than a complete technical null (Figure~\ref{fig:p1f6}E; Table~\ref{tab:p1rawcontrol}). In the anatomical permutation control, the observed p value was 0.0015 and 5.0\% of permuted comparisons fell below 0.05 (Figure~\ref{fig:p1f6}F). That fraction describes this permutation reference; it is not an independent repeated-dataset test of pipeline-wide type-I error.

The historical specimen-level scoring procedure used the pairwise AUC statistic in Equation~\eqref{eq:p1auc} and gave point AUCs of 0.822 for fibroblast-associated topic, 0.844 for the module score, 0.533 for PCA and 0.878 for NMF (Figure~\ref{fig:p1f7}A\label{call:p1f7}). No pairwise fixed-score bootstrap interval for an AUC difference excluded zero. The module score alone separated its procedure-specific permutation null at $p=0.0424$; fibroblast-associated topic and NMF gave $p=0.0604$ and $p=0.0597$, while PCA gave $p=0.4527$. Null means also differed from 0.5 (fibroblast-associated topic, 0.562; NMF, 0.553). These historical outputs pool raw coordinates from different fitted folds and include repeated patients; a permutation calculation does not make those coordinates comparable. They document the earlier procedure rather than supply evidence for a representation ranking (Tables~\ref{tab:p1sampleauc}, \ref{tab:p1correlations} and~\ref{tab:p1aucdifferences}).

The sample-level readouts had pairwise Spearman correlations from 0.174 to 0.710 (Figure~\ref{fig:p1f7}B). Their relatedness does not determine the ordering of their discrimination estimates.

\paperfigure{figure7_representation_inference}{fig:p1f7}{Historical specimen-unit scoring diagnostics in 14 DFU specimens (9 healing, 5 non-healing). (A) Points are leave-one-sample-out areas under the ROC curve (AUCs); horizontal bars are percentile 95\% intervals from 4,000 arm-stratified bootstrap draws of fixed out-of-fold scores, with common resampling indices across methods. Orange ticks mark the procedure-specific permutation-null means after rerunning outcome-dependent fitting and orientation. The dashed line at AUC=0.5 is a visual reference, not the fitted permutation null. Fibroblast topic denotes the frozen all-cell mean, Module the specified fibro-inflammatory score, PCA principal component analysis, and NMF nonnegative matrix factorization. Raw PCA/NMF coordinates can differ across fitted folds, so neither their pooled AUCs nor the displayed bootstrap contrasts support a representation ranking. (B) Cells contain Spearman correlations of these 14 specimen-level readouts; the colour scale and printed values encode the same statistic. Repeated patients and cross-fold comparability limit these historical diagnostics.}


\subsection{Same-model patient-pair discrimination and paired anatomical sensitivity}

The revised held-pair statistic in Equation~\eqref{eq:p1pairauc} evaluated all 28 healing/non-healing pairs from the 11 patients, fitting each scoring rule on the other nine patients. Within-model pair discrimination was 0.714 for the frozen topic (20 wins, no ties), 0.786 for the module (22 wins, no ties), 0.500 for PCA (14 wins, no ties) and 0.232 for NMF (five wins and three half-credit ties). These descriptive estimates use comparable scores within each pair and do not pool raw coordinates across fitted models (Table~\ref{tab:p1pairauc}; Figure~\ref{fig:p1f8}A\label{call:p1f8}). Changing held-patient expression or labels did not change fold preprocessing, component selection or orientation in the isolation tests; positive component rescaling also preserved NMF selection and pair order. The comparisons nevertheless reuse 11 discovery patients and a discovery-trained panel and encoder, and establish neither external validation nor a population representation ranking.

The earlier leave-one-patient-out procedure remains separately documented: topic AUC was 0.714 with add-one permutation $p=0.2991$, module AUC was 0.786 with $p=0.166$, PCA was 0.500 and NMF was 0.125. Its fixed-score bootstrap intervals and pairwise contrasts are historical procedure diagnostics, not uncertainty estimates for the revised held-pair estimand (Tables~\ref{tab:p1patientauc} and~\ref{tab:p1aucdifferences}). The change in the NMF estimate reflects the revised fold, projection and component-selection procedure; it is not an isolated estimate of a fold-scale effect.

\paperfigure{figure8_patient_unit_anatomy}{fig:p1f8}{Patient-unit discrimination and paired anatomical sensitivity. (A) Points are leave-one-patient-out AUCs for 7 healing and 4 non-healing patients after cell-count-weighted collapse of repeated specimens. Bars are 95\% intervals from 4,000 arm-stratified patient bootstrap draws of fixed out-of-fold scores; the dashed line marks AUC=0.5. Fibroblast topic, Module, PCA and NMF denote the frozen topic mean, fibro-inflammatory module, principal component and nonnegative matrix factorization readouts. (B) The point is the mean all-cell fibroblast-associated topic loading difference, foot minus forearm, across 10 paired patients; the bar is a 4,000-draw paired-patient bootstrap 95\% interval and the horizontal line is zero. The exhaustive sign-flip test with add-one correction gives p=0.0107 and Benjamini--Hochberg q=0.0402 across 15 topics. Both panels use the discovery source and are not independent clinical validation.}


The paired anatomical control confirmed that the frozen encoder responds to a known tissue difference. Across 10 paired patients, the fibroblast-associated topic increased in foot relative to forearm skin by +0.0476 (exhaustive sign-flip $p=0.0107$, BH $q=0.0402$); the fibroblast-restricted contrast was +0.1398 ($p=0.0146$). Four of 15 topics remained significant after BH correction (Figure~\ref{fig:p1f8}B; Table~\ref{tab:p1anatomy}). The earlier 3-of-15 sample-level positive-control result is directionally consistent, but the paired analysis is the more appropriate patient-unit sensitivity check. These controls make a complete assay failure less likely, but they do not turn the underpowered DFU comparison into a proof of no association.

\subsection{External projections remain construct audits}

The external immune-axis analysis used mutually exclusive lineage compartments. GSE223964 and GSE245703 gave immune-axis correlations of −0.024 and +0.846, respectively (Figure~\ref{fig:p1f9}A\label{call:p1f9}). Within the immune compartment, B/plasma-cell share was weakly related to the topic signal (−0.071 and +0.077), while correlations with median detected genes per cell were negative in both datasets (−0.429 and −0.427; Figure~\ref{fig:p1f9}B). The signs are inconsistent for a generalized immune-axis effect, and neither dataset has an authoritative patient-level healing endpoint (Table~\ref{tab:p1external}).

\input{figure_captions/f9}

A separate projection into GSE231643 provides a weaker, label-only concordance of direction: the within-fibroblast ratio falls from 2.89 in the discovery analysis to 1.66 (sample-bootstrap 95\% interval 1.01--6.10; two-sided Mann--Whitney $p=0.1429$; Table~\ref{tab:p1lineage}). The separate GSE241132 acute-wound cohort gave a broad immune-axis rho of −0.070 and an author-annotated B/plasma-cell rho of +0.210 across three donors; both values are descriptive and were not combined with the DFU outcome analyses. GSE268834 supplied another descriptive projection: the topic-associated fibroblast signal correlated with exclusive immune fraction at +0.714, and the diabetic-minus-non-diabetic fibroblast-topic contrast was +0.0464 with $p=0.2812$ (Figure~\ref{fig:p1f9}C). This is compatible with an immune/tissue-context hypothesis but cannot distinguish it from other between-sample differences.

GSE248247 was a deliberately low-yield stress test of the frozen projection. It contained 561 raw cells and 358 quality-controlled cells, with 6,994 of 7,002 panel genes covered in each sample. The fibroblast topic contrast between one diabetic and one non-diabetic sample was −0.0204 with descriptive permutation $p=1.0$ (Figure~\ref{fig:p1f9}C). All correlations were non-estimable because n=2. The exclusive compartment sum was exactly one in both samples, demonstrating that the lineage gate was non-overlapping; the dataset cannot validate healing or support pooled patient inference.
